\documentclass[../../main.tex]{subfiles}
\begin{document}
\section{多重线性代数}\label{section:ch12-多重线性代数}

我们已经看到余向量（向量空间上取实值的线性函数）在流形理论中扮演的某些重要角色。最简单地说，张量就是一个或多个变量上取实值的多重线性函数；简单的例子包括余向量、内积与行列式。为给张量研究搭好舞台，本节在一般框架下发展多重线性函数的一些基本性质。

设 $V_1,\dots,V_k$ 与 $W$ 为向量空间。称映射 $F\colon V_1\times\cdots\times V_k\to W$ 为\textbf{多重线性的}，若当其他变量固定时，它分别作为每个变量的函数都是线性的：对每个 $i$，
\begin{align}
F(v_1,\dots,a v_i+a'v_i',\dots,v_k)
=aF(v_1,\dots,v_i,\dots,v_k)+a'F(v_1,\dots,v_i',\dots,v_k).\label{align:ch12-多重线性}
\end{align}
（一个变量的多重线性函数就是线性函数；两个变量的多重线性函数通常称为双线性函数。）记 $L(V_1,\dots,V_k;W)$ 为从 $V_1\times\cdots\times V_k$ 到 $W$ 的所有多重线性映射的集合。它在逐点加法与数乘下构成向量空间：
\begin{align*}
(F+F')(v_1,\dots,v_k)&=F(v_1,\dots,v_k)+F'(v_1,\dots,v_k),\\
(aF)(v_1,\dots,v_k)&=a\big(F(v_1,\dots,v_k)\big).
\end{align*}

下面几个例子值得记住。

\begin{example}[一些熟悉的多重线性函数]\label{example:ch12-熟悉多重线性}
\begin{enumerate}[(a)]
\item $\mathbb{R}^n$ 中的点积是两个向量上取标量值的双线性函数，用于计算向量长度与它们之间的夹角。
\item $\mathbb{R}^3$ 中的叉积是两个向量上取向量值的双线性函数，用于计算平行四边形面积，并求与两个给定向量正交的第三个向量。
\item 行列式是 $\mathbb{R}^n$ 中 $n$ 个向量上取实值的多重线性函数，用于判别线性无关并计算这些向量张成的平行六面体体积。
\item Lie 代数 $\mathfrak g$ 中的括号是 $\mathfrak g$ 中两个元素上取 $\mathfrak g$ 值的双线性函数。
\end{enumerate}
\end{example}

下一个例子也许不那么熟悉，但极为重要。

\begin{example}[余向量的张量积]\label{example:ch12-余向量张量积}
设 $V$ 为向量空间，$\omega,\eta\in V^*$。定义函数 $\omega\otimes\eta\colon V\times V\to\mathbb{R}$ 为
\begin{align}
\omega\otimes\eta(v_1,v_2)=\omega(v_1)\eta(v_2),\label{align:ch12-余向量张量积}
\end{align}
右端乘积只是实数的普通乘法。$\omega$ 与 $\eta$ 的线性性保证 $\omega\otimes\eta$ 是 $v_1,v_2$ 的双线性函数，故它是 $L(V,V;\mathbb{R})$ 的元素。例如，若 $(\varepsilon^1,\varepsilon^2)$ 表示 $(\mathbb{R}^2)^*$ 的标准对偶基，则 $\varepsilon^1\otimes\varepsilon^2\colon\mathbb{R}^2\times\mathbb{R}^2\to\mathbb{R}$ 是双线性函数
\begin{align*}
\varepsilon^1\otimes\varepsilon^2((w,x),(y,z))=wy.
\end{align*}
\end{example}

上例可推广到任意取实值的多重线性函数如下：设 $V_1,\dots,V_k,W_1,\dots,W_l$ 为实向量空间，$F\in L(V_1,\dots,V_k;\mathbb{R})$，$G\in L(W_1,\dots,W_l;\mathbb{R})$。定义函数
\begin{align*}
F\otimes G\colon V_1\times\cdots\times V_k\times W_1\times\cdots\times W_l\to\mathbb{R}
\end{align*}
为
\begin{align}
F\otimes G(v_1,\dots,v_k,w_1,\dots,w_l)=F(v_1,\dots,v_k)G(w_1,\dots,w_l).\label{align:ch12-函数张量积}
\end{align}
由 $F$ 与 $G$ 的多重线性性，$F\otimes G(v_1,\dots,v_k,w_1,\dots,w_l)$ 分别线性依赖于每个变量 $v_i$ 或 $w_j$，故 $F\otimes G\in L(V_1,\dots,V_k,W_1,\dots,W_l;\mathbb{R})$，称为 $F$ 与 $G$ 的\textbf{张量积}。

\begin{exercise}\label{exercise:ch12-03}
证明张量积运算双线性且结合：$F\otimes G$ 双线性依赖于 $F$ 与 $G$，且 $(F\otimes G)\otimes H=F\otimes(G\otimes H)$。
\end{exercise}

由上一习题的结果，我们可以无歧义地写出三个或更多多重线性函数的张量积而不加括号。若 $F_1,\dots,F_l$ 分别依赖 $k_1,\dots,k_l$ 个变量，则它们的张量积 $F_1\otimes\cdots\otimes F_l$ 是 $k=k_1+\cdots+k_l$ 个变量的多重线性函数，其作用方式为把前 $k_1$ 个向量代入 $F_1$、接下来 $k_2$ 个代入 $F_2$，依此类推，再把结果相乘。例如，若 $F,G$ 是两个向量的多重线性函数，$H$ 是三个向量的，则
\begin{align*}
F\otimes G\otimes H(v_1,\dots,v_7)=F(v_1,v_2)G(v_3,v_4)H(v_5,v_6,v_7).
\end{align*}
若 $\omega^j\in V_j^*$，$j=1,\dots,k$，则 $\omega^1\otimes\cdots\otimes\omega^k\in L(V_1,\dots,V_k;\mathbb{R})$ 是由
\begin{align}
\omega^1\otimes\cdots\otimes\omega^k(v_1,\dots,v_k)=\omega^1(v_1)\cdots\omega^k(v_k)\label{align:ch12-多重张量积}
\end{align}
给出的多重线性函数。

张量积运算之所以重要，部分是因为下面命题。命题中因指标繁多记号略显难看，但基本想法很简单：任意多重线性函数空间的一组基，可由基余向量的所有可能张量积组成。

\begin{proposition}[多重线性函数空间的一组基]\label{proposition:ch12-多重线性函数基}
设 $V_1,\dots,V_k$ 为维数分别为 $n_1,\dots,n_k$ 的实向量空间。对每个 $j\in\{1,\dots,k\}$，设 $(E^{(j)}_1,\dots,E^{(j)}_{n_j})$ 是 $V_j$ 的一组基，$(\varepsilon^{(j)}_1,\dots,\varepsilon^{(j)}_{n_j})$ 是 $V_j^*$ 的相应对偶基。则集合
\begin{align*}
B=\left\{\varepsilon^{(1)}_{i_1}\otimes\cdots\otimes\varepsilon^{(k)}_{i_k}:1\leqslant i_1\leqslant n_1,\dots,1\leqslant i_k\leqslant n_k\right\}
\end{align*}
是 $L(V_1,\dots,V_k;\mathbb{R})$ 的一组基，故其维数等于 $n_1\cdots n_k$。
\end{proposition}

\begin{proof}
需证 $B$ 线性无关且张成 $L(V_1,\dots,V_k;\mathbb{R})$。任取 $F\in L(V_1,\dots,V_k;\mathbb{R})$。对每个满足 $1\leqslant i_j\leqslant n_j$ 的有序 $k$ 元组 $(i_1,\dots,i_k)$，定义数
\begin{align}
F_{i_1\dots i_k}=F(E^{(1)}_{i_1},\dots,E^{(k)}_{i_k}).\label{align:ch12-分量}
\end{align}
我们证明
\begin{align*}
F=F_{i_1\dots i_k}\,\varepsilon^{(1)}_{i_1}\otimes\cdots\otimes\varepsilon^{(k)}_{i_k}
\end{align*}
（如往常使用求和约定），由此 $B$ 张成 $L(V_1,\dots,V_k;\mathbb{R})$。对任意 $k$ 元向量组 $(v_1,\dots,v_k)\in V_1\times\cdots\times V_k$，写 $v_1=v^{i_1}_1 E^{(1)}_{i_1},\dots,v_k=v^{i_k}_k E^{(k)}_{i_k}$，计算
\begin{align*}
F_{i_1\dots i_k}\,\varepsilon^{(1)}_{i_1}\otimes\cdots\otimes\varepsilon^{(k)}_{i_k}(v_1,\dots,v_k)
=F_{i_1\dots i_k}\,\varepsilon^{(1)}_{i_1}(v_1)\cdots\varepsilon^{(k)}_{i_k}(v_k)
=F_{i_1\dots i_k}v^{i_1}_1\cdots v^{i_k}_k,
\end{align*}
而由多重线性性 $F(v_1,\dots,v_k)$ 等于同一物。这就证明了结论。

为证 $B$ 线性无关，设有线性组合为零：
\begin{align*}
F_{i_1\dots i_k}\,\varepsilon^{(1)}_{i_1}\otimes\cdots\otimes\varepsilon^{(k)}_{i_k}=0.
\end{align*}
将其作用于任意有序 $k$ 元基向量组 $(E^{(1)}_{j_1},\dots,E^{(k)}_{j_k})$。由与上面相同的计算，这蕴含每个系数 $F_{j_1\dots j_k}$ 为零。因此 $B$ 的元素唯一线性组合为零者是平凡组合。
\end{proof}

顺便指出，该证明表明多重线性函数 $F$ 在 $B$ 中基下的分量 $F_{i_1\dots i_k}$ 由\eqref{align:ch12-分量}给出。故 $F$ 完全由它在基向量所有可能序列上的作用决定。

\subsection*{向量空间的抽象张量积}

上一节结果表明多重线性函数空间 $L(V_1,\dots,V_k;\mathbb{R})$ 可视为形如 $\omega^1\otimes\cdots\otimes\omega^k$（$\omega^1,\dots,\omega^k$ 为余向量）的对象的所有线性组合之集。本节给出一个构造，使这种张量积的线性组合在更抽象的框架下也有意义。构造略繁，但想法简单：给定有限维向量空间 $V_1,\dots,V_k$，我们将构造一个新向量空间 $V_1\otimes\cdots\otimes V_k$，其维数为各 $V_i$ 维数之积，它由形如 $v_1\otimes\cdots\otimes v_k$（$v_i\in V_i$）的对象的"形式线性组合"组成，并使 $v_1\otimes\cdots\otimes v_k$ 分别线性依赖于每个 $v_i$。（本节引入的许多概念——至少那些不显式涉及有限基的部分——在无限维情形同样成立；但为简单起见我们主要限于有限维。）

首先需要弄清"形式线性组合"。设 $S$ 为一集合。粗略地说，$S$ 元素的形式线性组合是形如 $\sum_{i=1}^m a_i x_i$ 的表达式，其中 $a_1,\dots,a_m$ 为实数，$x_1,\dots,x_m\in S$。当然，由于我们不假设 $S$ 有任何代数结构，不能真正把 $S$ 中元素相加或数乘。但这种表达式的本质特征是：它完全由 $S$ 中哪些元素出现在和里、以及它们各带什么系数决定。于是定义：对任意集合 $S$，$S$ 元素的\textbf{形式线性组合}是一个函数 $f\colon S\to\mathbb{R}$，使得除有限多个 $s\in S$ 外 $f(s)=0$。$S$ 上的\textbf{自由（实）向量空间}，记为 $F(S)$，是 $S$ 元素所有形式线性组合之集。在逐点加法与数乘下，$F(S)$ 成为 $\mathbb{R}$ 上的向量空间。

对每个 $x\in S$，有函数 $\delta_x\in F(S)$，它在 $x$ 取 $1$、在 $S$ 的其他元素取 $0$；我们通常把这个函数等同于 $x$ 本身，从而把 $S$ 视为 $F(S)$ 的子集。于是每个 $f\in F(S)$ 可唯一写为 $f=\sum_{i=1}^m a_i x_i$，其中 $x_1,\dots,x_m$ 是使 $f(x_i)\neq 0$ 的 $S$ 元素，$a_i=f(x_i)$。故 $S$ 是 $F(S)$ 的一组基，从而 $F(S)$ 有限维当且仅当 $S$ 是有限集。

\begin{proposition}[自由向量空间的特征性质]\label{proposition:ch12-自由向量空间特征}
对任意集合 $S$ 与任意向量空间 $W$，每个映射 $A\colon S\to W$ 都唯一扩张为线性映射 $\widehat A\colon F(S)\to W$。
\end{proposition}

\begin{exercise}\label{exercise:ch12-06}
证明上一命题。
\end{exercise}

现设 $V_1,\dots,V_k$ 为实向量空间。先作自由向量空间 $F(V_1\times\cdots\times V_k)$，即所有 $k$ 元组 $(v_1,\dots,v_k)$（$v_i\in V_i$）的有限形式线性组合之集。令 $R$ 为 $F(V_1\times\cdots\times V_k)$ 中由下列形式元素张成的子空间：
\begin{align}
(v_1,\dots,a v_i,\dots,v_k)-a(v_1,\dots,v_i,\dots,v_k),\qquad
(v_1,\dots,v_i+v_i',\dots,v_k)-(v_1,\dots,v_i,\dots,v_k)-(v_1,\dots,v_i',\dots,v_k),\label{align:ch12-关系子空间}
\end{align}
其中 $v_j,v_j'\in V_j$，$i\in\{1,\dots,k\}$，$a\in\mathbb{R}$。

定义空间 $V_1,\dots,V_k$ 的\textbf{张量积}，记为 $V_1\otimes\cdots\otimes V_k$，为如下商向量空间：
\begin{align*}
V_1\otimes\cdots\otimes V_k=F(V_1\times\cdots\times V_k)/R,
\end{align*}
并令 $\pi\colon F(V_1\times\cdots\times V_k)\to V_1\otimes\cdots\otimes V_k$ 为自然投影。元素 $(v_1,\dots,v_k)$ 在 $V_1\otimes\cdots\otimes V_k$ 中的等价类记为
\begin{align}
v_1\otimes\cdots\otimes v_k=\pi(v_1,\dots,v_k),\label{align:ch12-抽象张量积}
\end{align}
称为 $v_1,\dots,v_k$ 的（抽象）张量积。由定义，抽象张量积满足
\begin{align*}
v_1\otimes\cdots\otimes a v_i\otimes\cdots\otimes v_k&=a(v_1\otimes\cdots\otimes v_i\otimes\cdots\otimes v_k),\\
v_1\otimes\cdots\otimes(v_i+v_i')\otimes\cdots\otimes v_k&=(v_1\otimes\cdots\otimes v_i\otimes\cdots\otimes v_k)+(v_1\otimes\cdots\otimes v_i'\otimes\cdots\otimes v_k).
\end{align*}

注意定义蕴含 $V_1\otimes\cdots\otimes V_k$ 的每个元素可表为形如 $v_1\otimes\cdots\otimes v_k$（$v_i\in V_i$）的元素的线性组合；但一般而言张量积空间的每个元素未必形如 $v_1\otimes\cdots\otimes v_k$（见习题~12-1）。

\begin{proposition}[张量积空间的特征性质]\label{proposition:ch12-张量积特征}
设 $V_1,\dots,V_k$ 为有限维实向量空间。若 $A\colon V_1\times\cdots\times V_k\to X$ 是到向量空间 $X$ 的任意多重线性映射，则存在唯一线性映射 $\widehat A\colon V_1\otimes\cdots\otimes V_k\to X$ 使下图交换：
\begin{align*}
\begin{array}{ccc}
V_1\times\cdots\times V_k & \xrightarrow{A} & X,\\
\pi\downarrow && \ne\widehat A\\
V_1\otimes\cdots\otimes V_k &&
\end{array}
\end{align*}
其中 $\pi(v_1,\dots,v_k)=v_1\otimes\cdots\otimes v_k$。
\end{proposition}

\begin{proof}
首先，由自由向量空间的特征性质，任意映射 $A\colon V_1\times\cdots\times V_k\to X$ 唯一扩张为线性映射 $\widehat A_0\colon F(V_1\times\cdots\times V_k)\to X$，其特征为 $\widehat A_0(v_1,\dots,v_k)=A(v_1,\dots,v_k)$。$A$ 多重线性恰好意味着子空间 $R$ 含于 $\widehat A_0$ 的核，因为
\begin{align*}
\widehat A_0(v_1,\dots,a v_i,\dots,v_k)=A(v_1,\dots,a v_i,\dots,v_k)=aA(v_1,\dots,v_i,\dots,v_k)=a\widehat A_0(v_1,\dots,v_i,\dots,v_k),
\end{align*}
对\eqref{align:ch12-关系子空间}中另一式有类似计算。因此 $\widehat A_0$ 下降为线性映射 $\widehat A\colon V_1\otimes\cdots\otimes V_k=F(V_1\times\cdots\times V_k)/R\to X$，满足 $\widehat A\circ\pi=\widehat A_0$。由于 $\pi$ 是包含 $V_1\times\cdots\times V_k\hookrightarrow F(V_1\times\cdots\times V_k)$ 再接投影，这蕴含 $\widehat A\circ\pi=A$。唯一性来自 $V_1\otimes\cdots\otimes V_k$ 的每个元素可表为 $v_1\otimes\cdots\otimes v_k$ 型元素的线性组合，而 $\widehat A$ 在这些元素上被 $\widehat A(v_1\otimes\cdots\otimes v_k)=\widehat A_0(v_1,\dots,v_k)=A(v_1,\dots,v_k)$ 唯一确定。
\end{proof}

称其为特征性质，是因为它在同构意义下唯一刻画了张量积（见习题~12-3）。下一个结果是命题~\ref{proposition:ch12-多重线性函数基}~对抽象张量积空间的类比。

\begin{proposition}[张量积空间的一组基]\label{proposition:ch12-张量积基}
设 $V_1,\dots,V_k$ 为维数分别为 $n_1,\dots,n_k$ 的实向量空间。对每个 $j=1,\dots,k$，设 $(E^{(j)}_1,\dots,E^{(j)}_{n_j})$ 是 $V_j$ 的一组基。则集合
\begin{align*}
C=\left\{E^{(1)}_{i_1}\otimes\cdots\otimes E^{(k)}_{i_k}:1\leqslant i_1\leqslant n_1,\dots,1\leqslant i_k\leqslant n_k\right\}
\end{align*}
是 $V_1\otimes\cdots\otimes V_k$ 的一组基，故其维数等于 $n_1\cdots n_k$。
\end{proposition}

\begin{proof}
形如 $v_1\otimes\cdots\otimes v_k$ 的元素按定义张成张量积空间；把这种表达式中每个 $v_i$ 用基展开，即知 $C$ 张成 $V_1\otimes\cdots\otimes V_k$。

为证 $C$ 线性无关，设 $C$ 中元素的某线性组合为零：
\begin{align*}
a_{i_1\dots i_k}\,E^{(1)}_{i_1}\otimes\cdots\otimes E^{(k)}_{i_k}=0.
\end{align*}
对每个有序指标 $k$ 元组 $(m_1,\dots,m_k)$，定义多重线性函数 $\theta_{m_1\dots m_k}\colon V_1\times\cdots\times V_k\to\mathbb{R}$ 为
\begin{align*}
\theta_{m_1\dots m_k}(v_1,\dots,v_k)=\varepsilon^{(1)}_{m_1}(v_1)\cdots\varepsilon^{(k)}_{m_k}(v_k),
\end{align*}
其中 $(\varepsilon^{(j)}_i)$ 是与 $(E^{(j)}_i)$ 对偶的 $V_j^*$ 的基。由于 $\theta_{m_1\dots m_k}$ 多重线性，由张量积的特征性质它下降为线性函数 $\widehat\theta_{m_1\dots m_k}\colon V_1\otimes\cdots\otimes V_k\to\mathbb{R}$。于是
\begin{align*}
0=\widehat\theta_{m_1\dots m_k}\left(a_{i_1\dots i_k}E^{(1)}_{i_1}\otimes\cdots\otimes E^{(k)}_{i_k}\right)
=a_{i_1\dots i_k}\theta_{m_1\dots m_k}(E^{(1)}_{i_1},\dots,E^{(k)}_{i_k})=a_{m_1\dots m_k},
\end{align*}
这表明 $C$ 线性无关。
\end{proof}

\begin{proposition}[张量积空间的结合性]\label{proposition:ch12-张量积结合}
设 $V_1,V_2,V_3$ 为有限维实向量空间。存在唯一同构
\begin{align*}
V_1\otimes(V_2\otimes V_3)\cong V_1\otimes V_2\otimes V_3\cong(V_1\otimes V_2)\otimes V_3,
\end{align*}
使得 $v_1\otimes(v_2\otimes v_3)$、$v_1\otimes v_2\otimes v_3$ 与 $(v_1\otimes v_2)\otimes v_3$ 型元素互相对应。
\end{proposition}

\begin{proof}
我们构造同构 $V_1\otimes V_2\otimes V_3\cong(V_1\otimes V_2)\otimes V_3$；另一个类似构造。由 $\alpha(v_1,v_2,v_3)=(v_1\otimes v_2)\otimes v_3$ 定义的映射 $\alpha\colon V_1\times V_2\times V_3\to(V_1\otimes V_2)\otimes V_3$ 显然多重线性，故由张量积的特征性质它唯一下降为线性映射 $\widehat\alpha\colon V_1\otimes V_2\otimes V_3\to(V_1\otimes V_2)\otimes V_3$，满足对所有 $v_1\in V_1,v_2\in V_2,v_3\in V_3$ 有 $\widehat\alpha(v_1\otimes v_2\otimes v_3)=(v_1\otimes v_2)\otimes v_3$。由于 $(V_1\otimes V_2)\otimes V_3$ 由 $(v_1\otimes v_2)\otimes v_3$ 型元素张成，$\widehat\alpha$ 满射，故由维数原因它是同构。它显然是唯一这种同构，因为任何其他同构都必须在张成 $V_1\otimes V_2\otimes V_3$ 的 $v_1\otimes v_2\otimes v_3$ 型元素集合上与 $\widehat\alpha$ 一致。
\end{proof}

抽象框架中的张量积与我们早先定义的多重线性泛函的更具体张量积之间的联系，基于下面的命题。

\begin{proposition}[抽象张量积与具体张量积]\label{proposition:ch12-抽象与具体张量积}
若 $V_1,\dots,V_k$ 为有限维向量空间，则存在典范同构
\begin{align}
V_1^*\otimes\cdots\otimes V_k^*\cong L(V_1,\dots,V_k;\mathbb{R}),\label{align:ch12-抽象具体同构}
\end{align}
在此同构下，由\eqref{align:ch12-抽象张量积}定义的抽象张量积对应于由\eqref{align:ch12-多重张量积}定义的余向量张量积。
\end{proposition}

\begin{proof}
先定义映射 $\Phi\colon V_1^*\times\cdots\times V_k^*\to L(V_1,\dots,V_k;\mathbb{R})$ 为
\begin{align*}
\Phi(\omega^1,\dots,\omega^k)(v_1,\dots,v_k)=\omega^1(v_1)\cdots\omega^k(v_k).
\end{align*}
（此式无求和约定。）右端分别线性依赖于每个 $v_i$，故 $\Phi(\omega^1,\dots,\omega^k)$ 确实是 $L(V_1,\dots,V_k;\mathbb{R})$ 的元素。容易验证 $\Phi$ 作为 $\omega^1,\dots,\omega^k$ 的函数是多重线性的，故由特征性质它唯一下降为线性映射 $\widehat\Phi\colon V_1^*\otimes\cdots\otimes V_k^*\to L(V_1,\dots,V_k;\mathbb{R})$，满足
\begin{align*}
\widehat\Phi(\omega^1\otimes\cdots\otimes\omega^k)(v_1,\dots,v_k)=\omega^1(v_1)\cdots\omega^k(v_k).
\end{align*}
由定义立即得 $\widehat\Phi$ 把抽象张量积映为余向量的张量积。它还把命题~\ref{proposition:ch12-张量积基}~给出的 $V_1^*\otimes\cdots\otimes V_k^*$ 的基映为命题~\ref{proposition:ch12-多重线性函数基}~给出的 $L(V_1,\dots,V_k;\mathbb{R})$ 的基，故它是同构。（尽管我们用了基来证 $\widehat\Phi$ 是同构，$\widehat\Phi$ 本身是典范定义的，不依赖任何基。）
\end{proof}

利用这个典范同构，此后我们用记号 $V_1^*\otimes\cdots\otimes V_k^*$ 既表示抽象张量积空间，也表示 $L(V_1,\dots,V_k;\mathbb{R})$，视问题方便而定。由于假设向量空间均有限维，我们还可把每个 $V_j$ 等同于其第二对偶 $V_j^{**}$，从而得到另一个典范等同
\begin{align*}
V_1\otimes\cdots\otimes V_k\cong L(V_1^*,\dots,V_k^*;\mathbb{R}).
\end{align*}

\subsection*{向量空间上的协变与反变张量}

设 $V$ 为有限维实向量空间。若 $k$ 为正整数，$V$ 上的一个\textbf{协变 $k$ 张量}是 $k$ 重张量积 $V^*\otimes\cdots\otimes V^*$ 的元素，我们通常把它视为 $V$ 的 $k$ 个元素上取实值的多重线性函数：
\begin{align*}
\alpha\colon \underbrace{V\times\cdots\times V}_{k\text{ 个}}\to\mathbb{R}.
\end{align*}
数 $k$ 称为 $\alpha$ 的\textbf{阶}。按约定，$0$ 张量就是一个实数。我们把 $V$ 上所有协变 $k$ 张量的向量空间简记为
\begin{align*}
T^k(V^*)=\underbrace{V^*\otimes\cdots\otimes V^*}_{k\text{ 个}}.
\end{align*}

\begin{example}[协变张量]\label{example:ch12-协变张量}
设 $V$ 为有限维向量空间。
\begin{enumerate}[(a)]
\item 每个线性泛函 $\omega\colon V\to\mathbb{R}$ 都是多重线性的，故协变 $1$ 张量就是余向量。于是 $T^1(V^*)=V^*$。
\item $V$ 上的协变 $2$ 张量是两个向量上取实值的双线性函数，亦称双线性型。一例是 $\mathbb{R}^n$ 上的点积；更一般地，每个内积都是协变 $2$ 张量。
\item 行列式作为 $n$ 个向量的函数，是 $\mathbb{R}^n$ 上的协变 $n$ 张量。
\end{enumerate}
\end{example}

为某些目的，需把协变张量概念推广如下。对任意有限维实向量空间 $V$，定义 $V$ 上阶为 $k$ 的\textbf{反变张量}空间为向量空间
\begin{align*}
T^k(V)=\underbrace{V\otimes\cdots\otimes V}_{k\text{ 个}}.
\end{align*}
特别地 $T^1(V)=V$，按约定 $T^0(V)=\mathbb{R}$。由于 $V$ 有限维，此空间可等同于 $k$ 个余向量上的多重线性泛函之集。

更一般地，对任意非负整数 $k,l$，定义 $V$ 上型为 $(k,l)$ 的\textbf{混合张量}空间为
\begin{align*}
T^{(k,l)}(V)=\underbrace{V\otimes\cdots\otimes V}_{k\text{ 个}}\otimes\underbrace{V^*\otimes\cdots\otimes V^*}_{l\text{ 个}}.
\end{align*}

\begin{corollary}[张量空间的基]\label{corollary:ch12-张量空间基}
设 $V$ 为 $n$ 维实向量空间。设 $(E_i)$ 是 $V$ 的任意一组基，$(\varepsilon^j)$ 是 $V^*$ 的对偶基。则下列集合分别构成 $V$ 上各张量空间的基：
\begin{align*}
\{\varepsilon^{i_1}\otimes\cdots\otimes\varepsilon^{i_k}\}&\quad\text{为 }T^k(V^*)\text{ 的基},\\
\{E_{i_1}\otimes\cdots\otimes E_{i_k}\}&\quad\text{为 }T^k(V)\text{ 的基},\\
\{E_{i_1}\otimes\cdots\otimes E_{i_k}\otimes\varepsilon^{j_1}\otimes\cdots\otimes\varepsilon^{j_l}\}&\quad\text{为 }T^{(k,l)}(V)\text{ 的基}.
\end{align*}
因此 $\dim T^k(V^*)=\dim T^k(V)=n^k$，$\dim T^{(k,l)}(V)=n^{k+l}$。
\end{corollary}

特别地，一旦为 $V$ 选定一组基，每个协变 $k$ 张量 $\alpha\in T^k(V^*)$ 可唯一写为
\begin{align*}
\alpha=\alpha_{i_1\dots i_k}\,\varepsilon^{i_1}\otimes\cdots\otimes\varepsilon^{i_k},
\end{align*}
其中 $n^k$ 个系数 $\alpha_{i_1\dots i_k}$ 由 $\alpha_{i_1\dots i_k}=\alpha(E_{i_1},\dots,E_{i_k})$ 决定。
\end{document}
